% Generated draft: requires mathematical and source review.
\documentclass[11pt]{article}
\usepackage[margin=25mm]{geometry}
\usepackage{fontspec}
\setmainfont{Latin Modern Roman}
\usepackage{amsmath,amssymb,mathtools,unicode-math}
\setmathfont{Latin Modern Math}
\usepackage{enumitem,needspace,xurl,hyperref}
\setlength{\parindent}{0pt}
\setlength{\parskip}{5pt}
\setlength{\emergencystretch}{3em}
\newcommand{\N}{\mathbb{N}}
\newcommand{\Z}{\mathbb{Z}}
\newcommand{\Q}{\mathbb{Q}}
\newcommand{\R}{\mathbb{R}}
\newcommand{\C}{\mathbb{C}}
\newcommand{\F}{\mathbb{F}}
\newcommand{\A}{\mathbb{A}}
\newcommand{\PP}{\mathbb{P}}
\newcommand{\E}{\mathbb{E}}
\newcommand{\eps}{\varepsilon}
\newcommand{\dd}{\,\mathrm{d}}
\begin{document}
\section*{UnsolvedMath ID: 20000859}
\textbf{Source problem code:} catalogue entry
\subsection{Definitions and mathematical statement}
\textbf{Chapter B question (Alessandra Angelucci).} Compare the spatial scales of classical receptive fields and surround modulation with horizontal, feedforward and feedback connections in primate visual cortex. The section describes tracer measurements related to receptive-field measurements, whose definitions depend on the experimental technique.

\textbf{Source observations and qualifications.} The notes distinguish minimum response field (mRF), high-contrast summation field (hsRF) and low-contrast summation field, but duplicate the abbreviation hsRF for the last two. The aid says this duplication occurs in the official PDF itself and the intended third measurement is low contrast, as established by an unidentified primary paper. The notes report horizontal spread commensurate with low-contrast V1 summation fields, with possible roles in expansion and collinear facilitation, and feedback spread commensurate with the full center-surround range, with possible contextual/global-to-local roles. These are reported comparisons and mechanism suggestions, not causal identifications.

The talk summary additionally reports feedback about ten times faster than horizontal connections and as fast as feedforward connections; patchy and orientation-specific feedback rather than diffuse nonspecific feedback; excitatory origins of both horizontal and feedback axons; and approximately $80\%$ excitatory versus $20\%$ inhibitory targets, expressly correcting a questioner's $98\%$ recollection. It also quotes ``no long-range inhibitory connections'' in this visual-cortical discussion. These empirical statements are retained as source claims with their scope unresolved, not universal laws. The V1-to-LGN discussion explicitly says no data were available, and the list of other cortical levels is incomplete.

\textbf{Reported direct-route exclusion certificate.} For a specified direct monosynaptic horizontal mechanism, let $d_-$ be a lower cortical-distance bound obtained from a declared visual-to-cortical mapping, $R_H^+$ an upper one-way reach, $v_H^+>0$ an upper propagation speed, and $t_{\mathrm{src}}^-$ and $\tau_H^-$ lower bounds on source activation time and fixed delay. Let observed onset belong to $[t_{\mathrm{obs}}^-,t_{\mathrm{obs}}^+]$. The aid reports exclusion if either
\[
 d_->R_H^+,\qquad\text{or}\qquad
 t_{\mathrm{obs}}^+<t_{\mathrm{src}}^-+\tau_H^-+\frac{d_-}{v_H^+}.
\]
All bounds favor the candidate route. Equality is not excluded. Spread diameters must be converted to one-way reach and cortical distances must carry the retinotopic uncertainty; speed ratios alone do not determine arrival latency. Missing the reach-truncated latency cone rules out only the specified direct route. Passing this necessary test does not establish that route, and excluding it does not prove feedback, since feedforward, recurrent, indirect and concurrent mechanisms remain possible.

\textbf{Reported radial scale nonidentifiability.} For a nonnegative radial profile $h(r)$ define disk response and normalized second-moment scale by
\[
 D_h(R)=2\pi\int_0^Rrh(r)\,\dd r,\qquad
 \sigma_h^2=\frac{\int_0^\infty r^3h(r)\,\dd r}{\int_0^\infty rh(r)\,\dd r}.
\]
Given finitely many sampled radii and a baseline $h_0\geq0$ with finite positive radial mass, finite third-weighted integral and a positive local lower bound on an open sampling annulus, the aid proves there are arbitrarily small nonnegative perturbations $h_+,h_-$ with identical responses at every sampled radius and identical total radial mass, but $\sigma_{h_+}^2\ne\sigma_{h_-}^2$. The perturbations can be supported in a subannulus containing no sampling boundary and have zero first-weighted but nonzero third-weighted perturbation integral. Thus finite disk-summation data and total mass do not identify this moment-based spatial scale, even locally around the qualifying baseline.

The aid describes contemporary causal literature as supporting feedback contributions while emphasizing condition-dependent concurrent routes, but supplies no full citations. Neither proposition converts anatomical/functional scale agreement into a uniquely identified circuit mechanism.
\subsection{Short English statement}
The section asks how anatomical connection ranges relate to visual receptive-field and surround scales. The aid supplies a favorable-bounds test that can exclude a specified direct horizontal route, plus a proof that finitely sampled radial responses do not identify a second-moment scale. Neither result proves feedback or makes the section's reported empirical comparisons universal.
\subsection{Sources}
\begin{itemize}
\item Local UnsolvedMath registry, permanent ID $20000859$, \texttt{catalogue entry}; canonical \texttt{aim-biology-notes.json}, \texttt{notes[23]}, tagged section.
\item AIM workshop ``Inference and prediction in neocortical circuits,'' Chapter B, ``Anatomical Substrates for Functional Responses of Neurons in the Primate Visual Cerebral Cortex,'' Alessandra Angelucci: \url{https://aimath.org/WWN/brain/brain.pdf}. The supplied verification identifies a genuine question within the section.
\item The aid refers to an unidentified primary paper for the low-contrast abbreviation and to current causal-feedback literature, without full references. No independent PDF/paper inspection is claimed here.
\item Research attempt 1, \texttt{partial\_result}, run status \texttt{partially\_solved}, internal rigor \texttt{proved}, low novelty confidence. Explicitly a machine-generated research aid, not peer reviewed; independently verify before citation. Original literature field empty.
\end{itemize}
\end{document}
